\section{Implikasi bagi Desain Bahasa dan Kompilator}
\label{sec:ch04-implikasi}

\subsection{Mengapa Leksikal Cukup Tipe 3}

Token memiliki struktur yang sederhana: identifier adalah urutan
huruf dan angka, bilangan adalah urutan digit, dan operator adalah
simbol tunggal atau ganda. Struktur seperti ini dapat dinyatakan
dengan ekspresi reguler dan dikenali dengan otomata hingga. Tidak ada
kebutuhan untuk menyimpan konteks bersarang pada level token.

Penggunaan otomata hingga memberi keuntungan efisiensi. Lexer dapat
berjalan dalam waktu linear terhadap panjang masukan dan menggunakan
memori konstan. Alat seperti Flex menghasilkan DFA yang sangat cepat
dari spesifikasi ekspresi reguler.

\subsection{Mengapa Sintaksis Memerlukan Tipe 2}

Struktur program memiliki sarang: kurung, blok, dan pemanggilan
fungsi. Bahasa seperti \(\{a^n b^n\}\) tidak regular karena
memerlukan penghitungan yang tidak terbatas. CFG dan otomata pushdown
menyediakan stack untuk menangani sarang tersebut.

Namun, CFG juga memiliki batas. CFG tidak dapat memeriksa bahwa
sebuah variabel telah dideklarasikan sebelum dipakai, atau bahwa tipe
operand cocok. Pemeriksaan seperti ini memerlukan informasi konteks
dan tabel simbol, sehingga berada di luar daya ekspresif CFG.

\subsection{Mengapa Semantik Melampaui Tipe 2}

Analisis semantik memerlukan atribut, lingkungan (\eng{environment}),
dan aturan tipe. Formalisme yang lebih sesuai adalah
\eng{attribute grammar} atau \eng{syntax-directed definition} (SDD).
Pada formalisme ini, setiap simpul pohon memiliki atribut yang
dihitung dengan aturan. Sebagian atribut bersifat
\eng{synthesized} (dihitung dari anak), sebagian
\eng{inherited} (diwarisi dari orang tua atau saudara).

Dengan attribute grammar, pemeriksaan tipe dapat dinyatakan secara
deklaratif. Sebagai contoh, atribut \code{type} pada simpul ekspresi
dihitung dari tipe operand. Jika operand tidak kompatibel, atribut
bernilai \code{error} dan diagnostik dilaporkan.

\subsection{Konsekuensi Desain}

\begin{table}[H]
\centering
\caption{Pemilihan formalisme untuk setiap fase}
\label{tab:ch04-formalisme}
\begin{tabularx}{\textwidth}{@{}l l X@{}}
\toprule
\textbf{Fase} & \textbf{Formalisme} & \textbf{Alasan} \\
\midrule
Leksikal & Regular/DFA & Struktur token sederhana, efisien \\
Sintaksis & CFG/PDA & Struktur bersarang memerlukan stack \\
Semantik & Attribute grammar & Memerlukan konteks dan tipe \\
Validasi & Aturan dan traversal & Memerlukan analisis lintas simpul \\
\bottomrule
\end{tabularx}
\end{table}

Pemilihan formalisme yang tepat memiliki konsekuensi praktis. Jika
sebuah bahasa dirancang dengan tata bahasa yang terlalu kompleks,
parser menjadi sulit dibangun dan dipelihara. Jika tata bahasa terlalu
sederhana, informasi penting harus diperiksa pada fase berikutnya.
Desain bahasa yang baik menempatkan setiap kekangan pada tingkat yang
paling sesuai.

\begin{warningbox}
Jangan menggunakan ekspresi reguler untuk memeriksa struktur
bersarang seperti kurung seimbang. Ekspresi reguler tidak memiliki
memori tak terbatas, sehingga akan gagal pada sarang yang dalam.
Gunakan parser atau PDA untuk keperluan tersebut.
\end{warningbox}
